% Zdroj: PDF priklady-podzim-2023.pdf, fyzická str. 18-19. Značka: specializace UI-SU. Zařazeno: S3.
\section{Kombinace více modelů}
\szzotazka{podzim 2023}{30}{Kombinace více modelů}

\noindent\textbf{Zadání.}\par
\begin{enumerate}
  \item Popište hlavní myšlenky technik bagging a boosting. Popište podrobněji jeden vybraný konkrétní algoritmus pro
    boosting.
  \item Jak funguje klasifikátor založený na náhodných lesech (random forest)? Jakým způsobem se trénuje?
  \item Porovnejte rozhodovací stromy a náhodné lesy. Jaké jsou výhody a nevýhody těchto modelů?
\end{enumerate}

\medskip
\noindent\textbf{Náčrt řešení.}\par
\begin{enumerate}
  \item Podstatou baggingu je vytvoření více modelů trénovaných na různých podmnožinách (bootstrap)\footnote{V originálním PDF
    je zde překlep: \uv{bootsrap}; zde uvedeno korektně.} trénovací množiny.
    Tyto podmnožiny jsou typicky samplovány s opakováním. Hlavním myšlenkou boostingu je trénování několika modelů
    postupně s tím, že další modely jsou trénovány na vážených datech, kde instance nesprávně klasifikované aktuálním
    klasifikátorem mají větší váhu. Příkladem může být např. algoritmus AdaBoost.
  \item Náhodné lesy obsahují množinu mnoha rozhodovacích stromů, které hlasují o celkovém výsledku klasifikace/regrese.
    Jednotlivé rozhodovací stromy jsou trénované na podmnožině instancí trénovací množiny a na podmnožině příznaků.
  \item Výhodou rozhodovacích stromů je jejich intrepretabilita a rychlost trénování. Výhodou náhodných lesů je jejich
    (typicky) větší přesnost a schopnost generalizace.
\end{enumerate}

\medskip
\noindent\textit{Doplňující vysvětlení (nad rámec oficiálního náčrtu):} Bootstrapový vzorek má stejnou velikost $N$ jako trénovací množina a~losuje se s~opakováním (asi $37\,\%$ příkladů v~něm chybí). Bagging snižuje hlavně rozptyl, boosting hlavně vychýlení. AdaBoost dá slabému klasifikátoru s~váženou chybou $\varepsilon_m$ hlas $\alpha_m=\tfrac12\ln\frac{1-\varepsilon_m}{\varepsilon_m}$ a~váhy chybně klasifikovaných příkladů vynásobí $e^{\alpha_m}$, správně klasifikovaných $e^{-\alpha_m}$ (pak normalizuje). V~náhodném lese se podmnožina příznaků losuje \emph{v~každém uzlu} znovu (typicky $\sqrt{D}$ příznaků) a~stromy se neprořezávají; to dekoreluje jejich chyby. Nevýhody lesa oproti jednomu stromu: ztráta interpretovatelnosti a~vyšší nároky na čas i~paměť při trénování i~predikci; nevýhodou jednoho stromu je vysoký rozptyl a~sklon k~přeučení.
